\subsection{Finding the Exponent} 
One of the things we most often want to do with an exponential model is to answer a question about the exponent.
\begin{Example}
The population of a small town increases by 3\% every year. If its population now is 12,000, how long will it take its population to reach 18,000? Round your answer to two decimal places.
\begin{lpic}{}{8}
\end{lpic}
\end{Example}
Sometimes---especially when only integers make sense---we want to understand our question as an inequality.
\begin{Example}
Alice deposits \$1500 in an account which pays 5.5\% interest, compounded annually. How long does it take for the value to be at least \$2500?
\begin{itemize}
\item Since the interest is compounded \makeblank[1in], interest is only added after a \makeblank passes.
\item The value ``is at least'' \$2500 means \makeblank.
\end{itemize}
\begin{lpic}{}{8}
\end{lpic}
\end{Example}